\usepackage{blindtext} %Alle benötigten Packages einbinden, wichtig um diese überhaupt nutzen zu können
\usepackage[utf8]{inputenc} % evtl. nötig
\usepackage{newclude}
\usepackage{graphicx}
\usepackage{adjustbox}
\usepackage{epigraph}
\usepackage[style=apa, doi=false, isbn = false]{biblatex} %Zitierstil IEEE, macht Literaturverzeichnis
\usepackage{caption}
\usepackage[ngerman]{babel} %Rechtschreibprüfung, aber nur rudimentär
\usepackage{csquotes}
\usepackage[T1]{fontenc}
\usepackage{newtxtext,newtxmath} % modernes Times New Roman ähnliches Paket
\usepackage{mathptmx}
\usepackage{multicol}
\usepackage{tabularx}
\usepackage{geometry} %Seitenränder etc. anpassbar machen
\usepackage{layout}
\usepackage{fancyhdr} %macht schöne Kopf- und Fußzeilen
\usepackage{moresize} %mehr Schriftgrößen möglich mit /Large /LARGE /HUGE /small etc.
\usepackage{setspace} %Zeilenabstand änderbar
\usepackage[colorlinks,pdfpagelabels,pdfstartview = FitH,bookmarksopen = true,bookmarksnumbered = true,linkcolor = black,plainpages = false,hypertexnames = false,citecolor = black,urlcolor = black]{hyperref} %Macht Verlinkungen im PDF, wie z.B. Kapitel im Inhaltsverzeichnis, Farben sind auf schwarz gesetzt, sind standardmäßig bunt
\usepackage[none]{hyphenat}
\usepackage{multirow}
\usepackage{tocloft}
\usepackage{minitoc}
\usepackage[page]{appendix}
\usepackage{etoolbox}
\usepackage{hyperref}
\usepackage{array}
\usepackage{longtable}
\usepackage{pdfpages}
\usepackage{booktabs}
\usepackage{array}
\usepackage{ragged2e} % for \RaggedRight
\usepackage{listings}
\usepackage{xcolor}
\usepackage{listings}
\usepackage{pdflscape}
\usepackage{rotating}
\usepackage{xcolor}

\lstset{
  language=R,
  basicstyle=\ttfamily\scriptsize,  % smaller font for A4
  breaklines=true,                   % wrap long lines
  breakatwhitespace=false,
  frame=single,                      % add frame around code
  backgroundcolor=\color[gray]{0.95}, % light gray background
  numbers=left,                      % line numbers
  numberstyle=\tiny,
  numbersep=5pt,
  showstringspaces=false,
  tabsize=2,
  columns=flexible,                  % better handling of proportional fonts
  literate=                           % preserve German Umlauts
    {Ö}{{\"O}}1
    {Ä}{{\"A}}1
    {Ü}{{\"U}}1
    {ß}{{\ss}}1
    {ö}{{\"o}}1
    {ä}{{\"a}}1
    {ü}{{\"u}}1,
  keywordstyle={},                   % remove keyword colors
  commentstyle= \color{darkgray}\itshape ,                   % remove comment colors
  stringstyle={}                     % remove string colors
}




\sloppy


% Define a new TOC file for the appendix
\newcommand{\appendixtoc}{appendix.toc}

% Command to add sections to the appendix TOC only
\newcommand{\appendixsection}[1]{%
  \section*{#1} % Section title without numbering
  \addtocontents{appendix}{\protect\contentsline{section}{#1}{\thepage}} % Add entry to appendix-specific TOC
}

% Set up a separate TOC for the appendix
\newcommand{\listofappendix}{%
  \section*{Anhang Inhaltsverzeichnis} % Title of the appendix TOC
  \addcontentsline{toc}{section}{Anhang Inhaltsverzeichnis} % Optional: add to main TOC
  \begin{itemize}
    \input{\appendixtoc} % Import appendix TOC entries
  \end{itemize}
}

\addbibresource{literatur.bib} %Literaturverzeichnis Datei einbinden

\geometry{ %Seitenränder festlegen
top=25mm,
bottom=20mm,
left= 25mm,
right=40mm,
headheight=15pt
}

\pagestyle{fancy} %Seitenstil auf fancyhdr festlegen
\fancypagestyle{plain}{} %Seitenstil plain wird bei Kapitelbeginn genutzt, soll auf fancy sein
\fancyhf{} %Resette Kopf- und Fußzeile
\fancyhead[L]{\nouppercase\leftmark} %Schreibe klein(Normal) Kapitel in die Kopfzeile links
\fancyhead[R]{\thepage} %Schreibe Seitenzahl in die Kopfzeile rechts